% P03.1 canonical-section persistent anchors proof, live integration.

\begin{proof}[\Cref{lem:p031-continuous-section-pointwise-openness}]
Let $N$ be a neighborhood of $\sigma(z)$.  Continuity of $\sigma$ makes
$\sigma^{-1}(N)$ a neighborhood of $z$, and the section identity gives
\[
                    \sigma^{-1}(N)\subset q(N).
\]
Thus $q(N)$ is a neighborhood of $z$.

For a neighborhood $U$ of $z$, the set $q^{-1}(U)$ is a neighborhood of
$\sigma(z)$, whence
\[
 (a\circ q)^*(\sigma(z))
   \leq \sup_{q^{-1}(U)}a\circ q
   =\sup_U a.
\]
Taking the infimum over $U$ gives the upper bound in
\eqref{eq:p031-continuous-section-usc-regularization}.  Conversely, for
every neighborhood $N$ of $\sigma(z)$, the set $q(N)$ is a neighborhood of
$z$ and
\[
                  \sup_N(a\circ q)=\sup_{q(N)}a\geq a^*(z).
\]
Taking the infimum over $N$ proves the reverse bound and hence the equality.
\end{proof}

\begin{proof}[\Cref{thm:p031-continuous-section-persistent-component-anchor}]
For each $i$, the set $\sigma(K_i)$ is connected, contains
$\sigma(\xi)$, and lies in $q^{-1}(K_i)$.  It is therefore contained in the
component $C_i$.  Applying $q$ gives
\[
        K_i=q(\sigma(K_i))\subset q(C_i)\subset K_i,
\]
which proves
\eqref{eq:p031-continuous-section-component-surjectivity}.

If $i\leq j$, then $C_j$ is a connected subset of $q^{-1}(K_i)$ containing
$\sigma(\xi)$, so $C_j\subset C_i$.  Finally $x\in K_i$ for every $i$, and
the first conclusion gives $\sigma(x)\in\sigma(K_i)\subset C_i$.  This proves
\eqref{eq:p031-continuous-section-component-nesting} and
\eqref{eq:p031-continuous-section-persistent-anchor}.
\end{proof}

\begin{proof}[\Cref{cor:p031-canonical-scalar-section-anchor}]
Poineau's Corollary~3.14 makes every point of $X$ universal.  By
Definition~3.2, its canonical lift to $X_{K'}$ is the point defined by the
multiplicative tensor norm on
$\mathcal H(x)\widehat\otimes_KK'$.  Corollary~3.7 makes the lift continuous
on every affinoid domain, and Remark~3.13 globalizes that conclusion to
analytic spaces.  The lift lies over $x$, so the resulting map
$\sigma_{K'/K}$ is a continuous section of $\pi_{K'/K}$.

Apply
\Cref{lem:p031-continuous-section-pointwise-openness} to obtain pointwise
openness and
\eqref{eq:p031-canonical-scalar-section-regularization}.  Apply
\Cref{thm:p031-continuous-section-persistent-component-anchor} to obtain the
remaining conclusions.
\end{proof}

\begin{proof}[\Cref{cor:p031-canonical-scalar-persistent-terminal-anchor}]
By the proof of
\Cref{thm:p031-rational-defect-terminal-boundary-core}, every $K_\alpha$ is
nonempty and connected, contains $x$ and $\xi$, and satisfies
\[
             K_\beta\subset K_\alpha
             \qquad(0<\alpha<\beta<g_v(x)).
\]
Apply
\Cref{thm:p031-continuous-section-persistent-component-anchor} to scalar
projection and its canonical continuous section from
\Cref{cor:p031-canonical-scalar-section-anchor}.  The theorem gives
\eqref{eq:p031-canonical-scalar-terminal-component}, the stated nesting, and
the explicit common anchor in
\eqref{eq:p031-canonical-scalar-persistent-terminal-anchor}.
\end{proof}
